% English Abstract

\begin{abstracten}

Large language model agents accomplish tasks through multi-step tool use in real environments such as code sandboxes and file systems, where reinforcement learning has become the dominant approach for continuous improvement. However, RL for interactive agents faces a fundamental bottleneck: both training data and trustworthy rewards are difficult to obtain at scale through human effort. Real user data accumulates slowly with narrow coverage; and rewards that rely solely on agent self-reports lead to reward hacking, where the model learns to falsely claim task completion.

This thesis designs and implements a self-evolving RL system for agentic LLMs, whose primary objective is to enable agents to autonomously generate training data, obtain trustworthy rewards, and produce next-step training tasks without human annotation. The system takes self-evolution as its first goal: beyond the initial seed, all subsequent training queries are autonomously generated based on prior execution experience—an observer summarizes the previous step's best trajectory, a questioner generates persona-consistent new seeds, and sandbox file-system state is inherited across steps so that the agent continues building on prior outputs. An oversample-eliminate-select mechanism further ensures that only the gradient-optimal trajectory subset is used for training at each step, improving self-evolution efficiency.

For autonomous data generation, three collaborative agents turn each sampling into a self-driven interaction: an actor solves tasks via ReAct in a sandbox, an observer computes deterministic state diffs, and a questioner generates next-step seeds from prior experience. For reward trustworthiness, a frozen judge scores based on real environment state diffs rather than agent self-reports, suppressing reward hacking at the evidence level. In addition, the system drives self-evolution of the evaluator and infrastructure from failure cases: it periodically collects failure cases, uses an LLM to attribute the root cause to either the model or the infrastructure, and accordingly revises prompts or distills missing capabilities into reusable skills, extending self-evolution from tasks and policy to the evaluator and infrastructure.

The system is validated end-to-end on Qwen3.5-9B, with the full self-evolution pipeline operational and sandbox state inheritance verified. Results show that a small number of seed queries can drive complete RL training, and the system operates without external data beyond the initial seed.

\end{abstracten}
